% Direct proof of the P06.1 algebraic-split toric-form extension.
% The splitting extension is exactly in the field class required by GJR26,
% Corollary 5.12; no assertion is made for a general analytic extension.

\begin{proof}[Proof of
\Cref{cor:p061-algebraic-split-toric-forms}]
We prove the product assertion, the assertion with no parameter being the
case where \(Y\) is a point.  Let \(\Omega\) be an ordinary open analytic
subspace of \(X^{\mathrm{an}}\times_KY\), let \(h\in\operatorname{PH}(\Omega)\),
and write
\[
 \Pi:X_L^{\mathrm{an}}\times_LY_L
       \longrightarrow X^{\mathrm{an}}\times_KY
\]
for scalar extension.  The space \(Y_L\) is good, strict, and quasi-smooth,
and \(\Pi^{-1}(\Omega)\) is an ordinary open analytic subspace of the
product of the split toric variety \(X_L^{\mathrm{an}}\) with \(Y_L\).
Scalar extension preserves classical pluriharmonicity by
\cite[Proposition~7.2(6)]{GR25Classical}.  Hence
\(\Pi^*h\) is classically pluriharmonic on \(\Pi^{-1}(\Omega)\), and
\Cref{cor:p061-split-toric-quasismooth-products} makes it
\(\mathbf R\)-PL there.

The ordinary open space \(\Omega\) is good and strict, and \(L\) is the
completion of an algebraic extension of \(K\).  Ground-field descent for
\(\mathbf R\)-PL functions, with coefficient group \(\mathbf R\), therefore
applies by \cite[Corollary~5.12]{GJR26} and gives
\(h\in\operatorname{PL}_{\mathbf R}(\Omega)\).  Both \(h\) and \(-h\) are
classically psh.  By \cite[Theorem~4.3]{GR25Classical}, both signs are
semipositive; hence \(h\) is Gubler--Jell--Rabinoff harmonic by
\cite[Definition~7.2]{GJR26}.

Finally, transfer a classically pluriharmonic function through the displayed
analytic local isomorphism and apply the product assertion on the resulting
ordinary open subspace.  The conclusions transfer back through that local
isomorphism.
\end{proof}
